% Part: normal-modal-logic
% Chapter: frame-correspondence
% Section: frames

\documentclass[../../../include/open-logic-section]{subfiles}

\begin{document}

\olfileid[pt]{nml}{frd}{fra}

\olsection{Frames}

\begin{defn}
  Um \emph{frame} é um par $\mModel{F} = \tuple{W,R}$, onde $W$ é um
  conjunto não vazio de mundos e $R$ é uma relação binária sobre~$W$. Um
  modelo $\mModel{M}$ é \emph{baseado em} um frame
  $\mModel{F} = \tuple{W,R}$ se e somente se
  $\mModel{M} = \tuple{W, R, V}$ para alguma valoração $V$.
\end{defn}

\begin{defn}
  Se $\mModel{F}$ é um frame, dizemos que $!A$ é \emph{válida em
    $\mModel{F}$,} $\mModel{F} \Entails !A$, se $\mSat{M}{!A}$ para todo
  modelo~$\mModel{M}$ baseado em~$\mModel{F}$.
  
  Se $\mClass{F}$ é uma classe de frames, dizemos que $!A$ é \emph{válida em
    $\mClass{F}$,} $\mClass{F} \Entails !A$, se e somente se
  $\mModel{F} \Entails !A$ para todo frame~$\mModel{F} \in \mClass{F}$.
\end{defn}

Os frames são interessantes porque a correspondência entre esquemas e
propriedades da relação de acessibilidade~$R$ ocorre no nível dos frames,
\emph{não no dos modelos}. Por exemplo, embora \Ax{T} seja verdadeira em
todos os modelos reflexivos, nem todo modelo no qual \Ax{T} é verdadeira é
reflexivo. Contudo, \emph{é} verdade não apenas que \Ax{T} é \emph{válida}
em todos os \emph{frames} reflexivos, mas também que todo frame no qual
\Ax{T} é válida é reflexivo.

\begin{rem}
A validade em uma classe de frames é um caso particular da noção de validade
em uma classe de modelos: $\mClass{F} \Entails !A$ se e somente se
$\mClass{C} \Entails !A$, onde $\mClass{C}$ é a classe de todos os modelos
baseados em um frame de~$\mClass{F}$.

Evidentemente, se !!a{formula} ou um esquema é válido, O.s.\@, válido com
respeito à classe de \emph{todos} os modelos, ele também é válido com respeito
a qualquer classe~$\mClass{F}$ de frames.
\end{rem}

\end{document}
